%% CV for UPMC Statistician Role - Zhengyi (Joe) Ou
%% Tailored to emphasize sleep/circadian research, statistical analysis, and R programming

\documentclass[11pt,a4paper,sans]{moderncv}
\moderncvstyle{banking}
\moderncvcolor{blue}

\renewcommand*{\firstnamestyle}[1]{{\fontsize{34}{36}\bfseries\upshape\color{color1}#1}}
\renewcommand*{\lastnamestyle}[1]{{\fontsize{34}{36}\bfseries\upshape\color{color1}#1}}
\renewcommand*{\sectionstyle}[1]{{\sectionfont\color{color1}#1}}

\usepackage[utf8]{inputenc}
\usepackage{hyperref}
\hypersetup{
    colorlinks=true,
    linkcolor=blue,
    filecolor=magenta,
    urlcolor=blue,
    pdftitle={Zhengyi (Joe) Ou - CV},
    pdfpagemode=FullScreen,
}
\usepackage[scale=0.80]{geometry}
\usepackage{import}

\name{Zhengyi (Joe)}{Ou}
\address{Atlanta, GA}{}{}
\phone[mobile]{+1 404-926-6260}
\email{joeou2002@gmail.com}
\extrainfo{\href{https://linkedin.com/in/zhengyiou}{LinkedIn}, \href{https://github.com/2zOu2}{GitHub}}

\begin{document}

\makecvtitle

\vspace{6pt}
\small{Biostatistics graduate (MSPH, Emory University) with research experience analyzing longitudinal health outcomes, including sleep-wake patterns and circadian activity data. Proficient in R for statistical analysis, data visualization, and reproducible research workflows. Led analyses of 2,648-patient cohorts and supported clinical trial programming with SAS; experienced in translating complex datasets into clear reports and manuscript-ready findings for peer review.}

\section{Core Competencies}
\vspace{1pt}
\begin{itemize}
\item \textbf{Statistical Analysis \& R Programming}: Longitudinal analysis, mixed-effects models, survival analysis, Bayesian methods, data visualization, R Shiny dashboards
\item \textbf{Research Data Management}: Data preparation, validation, quality control, analysis-ready dataset creation from clinical and real-world data sources
\item \textbf{Scientific Reporting}: Reproducible workflows, publication-ready tables and figures, statistical summaries for peer review, oral and written presentation
\item \textbf{Clinical Trial \& Regulatory}: SAP tableshells, TLF generation, ICH E9(R1), CDISC/SDTM/ADaM fundamentals, analysis population definitions
\item \textbf{Statistical Software}: R (primary), SAS, SQL/MySQL, Python, Git/GitHub, R Shiny
\end{itemize}

\section{Professional Experience}
\vspace{3pt}
\begin{itemize}

\item{\cventry{Oct 2025--Present}{Research Assistant}{Rollins School of Public Health, Emory University}{Atlanta, GA}{}{\vspace{1pt}
\begin{itemize}
    \item Analyzed actigraphy data from post-stroke patients using Bayesian Weibull proportional hazards models in R to characterize sleep-wake pattern recovery trajectories, as part of an NIH-funded study
    \item Built an accompanying R Shiny dashboard to help visualize recovery trajectories for clinical interpretation
\end{itemize}}}

\vspace{3pt}

\item{\cventry{Aug 2025--Dec 2025}{Clinical Research Data Analyst}{Emory University School of Medicine}{Atlanta, GA}{}{\vspace{1pt}
\begin{itemize}
    \item Led real-world clinical outcomes analyses for a longitudinal cohort of 2,648 patients and 3,452 knees evaluating comparative effectiveness of three orthobiologic treatments for knee osteoarthritis
    \item Built analysis-ready longitudinal datasets from EHR-derived clinical outcomes, including follow-up visit alignment, missingness assessment, and QC workflows
    \item Developed reproducible R workflows for data preprocessing, validation, statistical analysis, and generation of publication-ready tables and figures
    \item Performed adjusted linear mixed-effects modeling to evaluate treatment effectiveness across repeated follow-up visits over a 24-month observation period
\end{itemize}}}

\vspace{3pt}

\item{\cventry{May 2025--Aug 2025}{Biostatistician Intern}{Canming Data (CRO)}{Beijing, China}{}{\vspace{1pt}
\begin{itemize}
    \item Supported SAP tableshell development and statistical programming for a randomized, double-blind, placebo-controlled Phase III COPD clinical trial
    \item Prepared TLF shells and programming specifications for subject disposition, baseline characteristics, efficacy analyses, safety summaries, and adverse events
    \item Supported MMRM-based primary efficacy analysis for Week 12 change from baseline in trough FEV1, including LS mean estimates, treatment comparisons, and sensitivity analyses
    \item Worked with raw eCRF datasets and reviewed analysis population definitions (FAS, PPS, Safety Set, PK Set); supported QC checks for treatment groups, missing values, and output validation
\end{itemize}}}

\vspace{3pt}

\item{\cventry{Sep 2024--Dec 2025}{Graduate Teaching Assistant}{Emory University}{Atlanta, GA}{}{\vspace{1pt}
\begin{itemize}
    \item QTM 151 - Scientific Programming (Python); QTM 210 - Probability
\end{itemize}}}

\vspace{3pt}

\item{\cventry{2023--2024}{Research Assistant}{Rollins School of Public Health, Emory University}{Atlanta, GA}{}{\vspace{1pt}
\begin{itemize}
    \item Developed adaptive Thompson Sampling algorithms in Python for ERP-based brain-computer interface research; integrated NLP-based language models to improve real-time classification in clinical signal environments
\end{itemize}}}

\end{itemize}

\section{Education}
\vspace{1pt}
\begin{itemize}

\item{\cventry{2024--2026}{Master of Science in Public Health (MSPH), Biostatistics}{Emory University, Rollins School of Public Health}{Atlanta, GA}{}{}}

\vspace{3pt}

\item{\cventry{2020--2024}{Bachelor of Science (B.S.), Applied Mathematics and Statistics}{Emory University, College of Arts and Sciences}{Atlanta, GA}{}{}}

\end{itemize}

\section{Languages}
\vspace{1pt}
\begin{itemize}
\item Mandarin (native), English (fluent).
\end{itemize}

\section{Publications \& Presentations}
\vspace{1pt}
\begin{itemize}
\item Co-author. ``Comparative Efficacy of Cellular Injectates for Knee Osteoarthritis: A Retrospective Longitudinal Analysis of Real-World Patient Outcomes.'' Under review, American Journal of Sports Medicine, 2026.
\item Presenter. ``Adaptive Stimulus Selection via Thompson Sampling in ERP-Based Brain-Computer Interfaces.'' Society for Neuroscience Annual Meeting, 2023.
\end{itemize}

\section{References}
\vspace{1pt}
\begin{itemize}
\item{Available upon request.}
\end{itemize}

\end{document}
